\documentclass[10pt,a4paper,ssfamily]{exam}
\usepackage[brazil]{babel}
\usepackage[numbers,sort&compress]{natbib}
\usepackage[utf8]{inputenc}
\renewcommand{\baselinestretch}{1.0}
\usepackage{epsfig}
\usepackage{mhchem}
\usepackage{graphicx}
\usepackage{svg}
\usepackage{makeidx}
\usepackage{multicol}
\usepackage{multirow}
\usepackage{amssymb}
\usepackage{amsmath}
\DeclareMathOperator{\sen}{sen}
\newcommand{\listfont}{\bf\fontencoding{OT1}\sffamily\large}
\newcommand{\titlesfont}{\fontencoding{OT1}\sffamily\large}
\renewcommand{\baselinestretch}{1.0}
\renewcommand{\t}{\textrm}
\newcommand{\cc}[1]{[\textrm{#1}]}
\usepackage{enumitem}
\setlist{nosep, topsep=2pt, partopsep=0pt, parsep=0pt, itemsep=2pt}
\newcommand{\nomera}{
  \noindent
  \begin{tabular}{|p{0.7\textwidth}|p{0.3\textwidth}|}
  Nome: & RA: \\
  \hline
  \end{tabular}
}

\newcommand{\qini}{\begin{enumerate}[resume]\item}
\newcommand{\qend}{\end{enumerate}}

\newcommand{\oC}{$^\circ$C}
\newcommand{\e}[1]{$\times 10^{#1}$}

\begin{document}
\sffamily
\pagestyle{empty}
\nomera

\begin{center}
{\bf QG101 - Química Geral}\\
{\bf Bloco 9 - Forças Intermoleculares e Estados da Matéria}\\
\underline{Justifique suas respostas.}
\end{center}

\vspace{0.2cm}
\begin{center}{\bf Questões}\end{center}

% ---------------------------------------------------------------
\qini
Coloque as substâncias \ce{CH4}, \ce{H2O}, \ce{NH3} e \ce{Ne} em ordem
\textbf{crescente} de ponto de ebulição e indique, para cada uma, o
principal tipo de força intermolecular envolvido.

\vspace{0.5cm}
\qend

% ---------------------------------------------------------------
\qini
A tabela abaixo apresenta os pontos de ebulição dos hidretos do grupo 16:

\begin{center}
\begin{tabular}{lcccc}
\hline
 & \ce{H2S} & \ce{H2Se} & \ce{H2Te} & \ce{H2O} \\
\hline
P.E. (\oC) & $-60$ & $-41$ & $-2$ & $100$ \\
\hline
\end{tabular}
\end{center}

\begin{enumerate}
  \item Explique por que o ponto de ebulição cresce de \ce{H2S} para
        \ce{H2Te}.
  \item Se a água seguisse a mesma tendência (apenas forças de London),
        seu ponto de ebulição seria de aproximadamente $-90$\oC. Explique
        por que o valor real é tão maior.
\end{enumerate}

\vspace{0.3cm}
\qend

% ---------------------------------------------------------------
\qini
Para cada uma das espécies abaixo, identifique o \textbf{principal} tipo
de força intermolecular presente no estado líquido/sólido:
\begin{enumerate}
  \item \ce{HF}
  \item \ce{CO2}
  \item \ce{CH4}
  \item \ce{NH3}
  \item \ce{I2}
\end{enumerate}

\vspace{0.3cm}
\qend

% ---------------------------------------------------------------
\qini
Explique, em termos de estrutura e de ligações de hidrogênio, por que o
gelo flutua na água líquida. Por que essa propriedade é considerada
anômala em relação à maioria das substâncias?

\vspace{0.8cm}
\qend

% ---------------------------------------------------------------
\qini
O etanol (\ce{C2H5OH}) é completamente solúvel em água, mas o hexano
(\ce{C6H14}) praticamente não se dissolve. Explique essa diferença usando
o princípio "semelhante dissolve semelhante" e os tipos de força
intermolecular envolvidos em cada caso.

\vspace{0.8cm}
\qend

% ---------------------------------------------------------------
\qini
O butano (\ce{C4H10}, $M = 58$~g/mol) ferve a $-1$\oC, enquanto a acetona
(\ce{C3H6O}, $M = 58$~g/mol) ferve a $56$\oC.
\begin{enumerate}
  \item As duas substâncias têm a mesma massa molar. O que explica a
        grande diferença nos pontos de ebulição?
  \item Quais forças intermoleculares atuam em cada uma?
\end{enumerate}

\vspace{0.3cm}
\qend

% ---------------------------------------------------------------
\qini
À temperatura ambiente, \ce{F2} e \ce{Cl2} são gases, \ce{Br2} é líquido e
\ce{I2} é sólido, embora todos sejam moléculas apolares formadas pelo
mesmo tipo de ligação covalente simples.
\begin{enumerate}
  \item Qual força intermolecular é responsável pelas diferenças de
        estado físico observadas?
  \item Relacione essa tendência ao número de elétrons e à
        polarizabilidade de cada molécula.
\end{enumerate}

\vspace{0.3cm}
\qend

% ---------------------------------------------------------------
\qini
A glicerina (com três grupos \ce{-OH} por molécula) é muito mais viscosa
que a água, que por sua vez é mais viscosa que a gasolina.
\begin{enumerate}
  \item Relacione a viscosidade de cada líquido com o tipo e a quantidade
        de forças intermoleculares presentes.
  \item Por que a água apresenta tensão superficial mais alta que a
        maioria dos solventes orgânicos comuns?
\end{enumerate}

\vspace{0.3cm}
\qend

% ---------------------------------------------------------------
\qini
Quando o \ce{NaCl} se dissolve em água, os íons \ce{Na+} e \ce{Cl-} ficam
rodeados por moléculas de água (hidratação).
\begin{enumerate}
  \item Que tipo de interação (distinta das três discutidas neste bloco)
        ocorre entre os íons e as moléculas de água?
  \item As moléculas de água se orientam de forma específica ao redor de
        cada íon. Indique qual extremidade da molécula de água
        (\ce{O} ou \ce{H}) fica voltada para o \ce{Na+} e qual fica
        voltada para o \ce{Cl-}, justificando.
\end{enumerate}

\vspace{0.3cm}
\qend

% ---------------------------------------------------------------
\qini
A figura abaixo mostra a curva de aquecimento de uma amostra de água pura
sob pressão constante, partindo do estado sólido.

\begin{center}
\includesvg[width=0.55\textwidth]{figuras/bloco09_fig2_curva_aquecimento}
\end{center}

\begin{enumerate}
  \item Por que a temperatura permanece constante durante os trechos de
        fusão e de ebulição, mesmo com fornecimento contínuo de calor?
  \item Para onde vai a energia fornecida durante esses trechos
        constantes, já que ela não aumenta a temperatura?
  \item O trecho de ebulição é mais longo (no eixo do calor) que o trecho
        de fusão. O que isso indica sobre as energias envolvidas em cada
        mudança de estado?
\end{enumerate}

\vspace{0.3cm}
\qend

\end{document}
